\documentclass[11pt,a4paper]{article}

% ── Packages ────────────────────────────────────────────────────────────────
\usepackage[utf8]{inputenc}
\usepackage[T1]{fontenc}
\usepackage[margin=1in]{geometry}
\usepackage{amsmath, amssymb, amsthm, mathtools, bm}
\usepackage{graphicx}
\usepackage{booktabs}
\usepackage{tabularx}
\usepackage{enumitem}
\usepackage{hyperref}
\usepackage{xcolor}
\usepackage{listings}
\usepackage{float}
\usepackage{algorithm}
\usepackage{algpseudocode}

\lstset{
  basicstyle=\ttfamily\footnotesize,
  breaklines=true,
  columns=fullflexible,
  keepspaces=true,
  frame=single,
  showstringspaces=false,
  numberstyle=\tiny\color{gray},
  commentstyle=\color{gray!70!black}\itshape,
}
\hypersetup{
  colorlinks=true,
  linkcolor=blue!50!black,
  citecolor=blue!50!black,
  urlcolor=blue!60!black,
}

\newcommand{\Utarget}{U_{\text{target}}}
\newcommand{\Ucirc}{U_{\text{circ}}}
\newcommand{\Real}{\mathbb{R}}
\newcommand{\Complex}{\mathbb{C}}
\newcommand{\Expect}{\mathbb{E}}
\newcommand{\E}{\mathrm{e}}

\title{\textbf{Benchmark Analysis and Convergence-Speed Study\\
on 2-/3-Qubit Brickwork Targets:\\
Diagnosis and a Rotosolve-Based Refinement Proposal}}
\author{Ayman Tarig}
\date{}

\begin{document}
\maketitle

\begin{abstract}
We benchmark the hybrid-action GQE pipeline on $50$ Haar-random brickwork
unitaries each at $n=2$ and $n=3$ qubits, comparing against Qiskit's
\texttt{optimization\_level=3} transpile. The agent achieves essentially
perfect fidelity in both cases ($\bar F \ge 0.9997$), and on $n=2$ matches
or beats Qiskit on every structural axis. On $n=3$ a striking pattern
emerges: GQE \emph{beats Qiskit on circuit depth} (median $\Delta D = -5$)
while \emph{losing on CNOT count} (median $\Delta C = +5$). The agent is
exchanging two-qubit gates for parallelism. We diagnose this trade-off as
a joint failure of the QASER reward base \emph{and} the inner-refinement
optimiser: L-BFGS does not converge tightly enough on the smallest
($\sim$Khaneja-Glaser-optimal) circuits, so the agent compensates by
sampling slightly bigger but easier-to-fit ones. We propose adopting
the Rotosolve closed-form coordinate update --- which is exact for our
process-fidelity objective --- as a configurable optimiser at both the
inner and post-training refinement stages, and discuss the expected
fidelity / CNOT / depth trade-offs and the implementation surface area.
\end{abstract}

% ════════════════════════════════════════════════════════════════════════════
\section{Benchmark setup and headline numbers}
% ════════════════════════════════════════════════════════════════════════════

We evaluate the trained policy on $N = 50$ independently sampled
\emph{brickwork-Haar} target unitaries per qubit count. Each target is the
product of alternating layers of Haar-random 2-qubit gates on adjacent
qubit pairs (even-offset layers $(0,1)(2,3)\ldots$ followed by odd-offset
layers $(1,2)(3,4)\ldots$, repeated $2n$ times). This distribution is a
strict generalisation of Haar-random for circuit-synthesis benchmarking:
the resulting unitary is approximately Haar at $\mathrm{poly}(n)$ depth
but its \emph{combinatorial origin} (a brickwork pattern of 2-qubit
unitaries) gives Qiskit's transpiler a fairer reference target than a
fully unstructured Haar matrix. Each circuit is trained from scratch with
the production \texttt{config.yml} (tiny GPT-2, 100 epochs, $1024$ rollouts
/ epoch, QASER reward with log-infidelity multiplier, post-training
refinement with $16$ multi-restarts $\times$ $100$ L-BFGS steps each).

Table~\ref{tab:headline} summarises the fifty-circuit averages.

\begin{table}[H]
\centering
\small
\begin{tabular}{lcccc}
\toprule
& \multicolumn{2}{c}{$n = 2$} & \multicolumn{2}{c}{$n = 3$}\\
\cmidrule(lr){2-3}\cmidrule(lr){4-5}
quantity                & GQE   & Qiskit & GQE   & Qiskit \\
\midrule
$\bar{F}$               & $0.99993$ & $1.00$ & $0.99977$ & $1.00$ \\
$\min F$                & $0.9983$  & ---    & $0.9986$  & ---    \\
fraction with $F \ge 0.99$ & $100\%$ & ---    & $100\%$ & --- \\
\midrule
mean CNOT count $\bar{C}$ & $\bm{2.92}$ & $3.00$ & $23.58$ & $\bm{19.00}$ \\
mean depth $\bar{D}$    & $\bm{14.20}$ & $15.00$ & $\bm{56.72}$ & $62.00$ \\
mean total gates $\bar{G}$ & $\bm{22.14}$ & $24.00$ & $\bm{95.38}$ & $96.00$ \\
\midrule
median $\Delta C$ (GQE$-$Qiskit) & $0$  & --- & $+5$ & ---\\
median $\Delta D$ (GQE$-$Qiskit) & $-1$ & --- & $-5$ & ---\\
median $\Delta G$ (GQE$-$Qiskit) & $-2$ & --- & $-1$ & ---\\
\midrule
$\bar{t}_{\text{train}}$ per circuit & \multicolumn{2}{c}{$75.5$ s} & \multicolumn{2}{c}{$120.2$ s}\\
\bottomrule
\end{tabular}
\caption{Per-instance averages over $50$ brickwork-Haar targets. Bold
indicates the better value. The Khaneja-Glaser CNOT lower bound is $3$
(for $n=2$) and $14$ (for $n=3$).}
\label{tab:headline}
\end{table}

The high-level diagnosis from these numbers is:

\begin{itemize}[leftmargin=2em]
\item \textbf{The agent has solved $n=2$.} It hits or beats Qiskit on
all three structural axes simultaneously, fidelity is essentially $1.00$,
and the search is robust ($5/50$ runs find $2$-CNOT decompositions for
non-maximally-entangling targets, while $44/50$ match the canonical
$3$-CNOT KAK form).

\item \textbf{$n = 3$ is good on fidelity and depth, bad on CNOT count.}
Forty-nine of fifty $3$-qubit instances use \emph{more} CNOTs than Qiskit
(Q1/median/Q3 deltas: $+3$/$+5$/$+6$), and forty-four of fifty have
\emph{shorter} depth than Qiskit (median $\Delta D = -5$). The total gate
count comes out roughly equal (median $\Delta G = -1$), so this is not
the agent ``inflating'' the circuit, it is the agent
\emph{redistributing} structure: trading sequential CNOTs for parallel
ones. Two-qubit gates are the dominant noise contributor on near-term
hardware, so this trade-off is not free.
\end{itemize}

% ════════════════════════════════════════════════════════════════════════════
\section{Diagnosis of the CNOT-vs-depth trade-off on $n=3$}
\label{sec:diagnosis}
% ════════════════════════════════════════════════════════════════════════════

Why does the agent prefer $\sim 24$ shallow CNOTs over $\sim 19$ deeper
ones, when the Pareto archive contains both options and the QASER reward
explicitly penalises CNOTs more heavily than rotations
($w_C = 2$, $w_G = 1.5$, $w_D = 1$)?

\subsection{The reward base does prefer the Qiskit-shaped circuit}

Let $M_D, M_C, M_G$ denote the running maxima of depth, CNOT count and
total-gate count seen so far in training. At the typical post-training
values $M_D \approx 80$, $M_C \approx 30$, $M_G \approx 120$ on $n=3$,
the QASER \emph{base} term is, for representative circuits,
\begin{align*}
\text{Qiskit-shaped (D=62, C=19, G=96):}\quad
& 1.0\cdot\tfrac{80}{63} + 2.0\cdot\tfrac{30}{20} + 1.5\cdot\tfrac{120}{97}
\;\approx\; 1.27 + 3.00 + 1.86 \;=\; 6.13,\\[2pt]
\text{GQE-shaped (D=57, C=24, G=95):}\quad
& 1.0\cdot\tfrac{80}{58} + 2.0\cdot\tfrac{30}{25} + 1.5\cdot\tfrac{120}{96}
\;\approx\; 1.38 + 2.40 + 1.88 \;=\; 5.66.
\end{align*}
The Qiskit-shaped (low-CNOT, taller) circuit produces a \emph{larger}
reward base by about $8\%$. So the reward formulation isn't the
problem. The agent simply isn't picking the Qiskit-shaped circuit.

\subsection{The bottleneck is angle convergence, not the reward}

The standard pattern in our 3-qubit training trajectories is:

\begin{enumerate}[leftmargin=2em]
\item Early epochs (1--20): the agent samples short circuits, gets
mediocre fidelity ($F \in [0.4, 0.7]$).
\item Middle epochs (20--60): the agent extends circuits to $\sim 100$
gates, finds layouts with $\sim 25$ CNOTs that yield post-inner-refinement
$F \approx 0.99$.
\item Late epochs (60--100): the policy concentrates around the
``$\sim 25$-CNOT, parallel'' regime. It rarely visits the
$\sim 19$-CNOT regime that Qiskit uses.
\end{enumerate}

The reason the agent never settles on the Qiskit-shaped regime is that
the inner-refinement L-BFGS, with $12$ steps and no multi-restart, fails
to converge \emph{cleanly} when the discrete circuit is close to the
Khaneja-Glaser CNOT lower bound. Tight, near-optimal-CNOT circuits live
in the most non-convex part of the variational loss surface (every
parameter contributes to many entries of $\Ucirc$ simultaneously, so the
condition number explodes). At the inner stage the optimiser ``gives up''
early and reports $F \approx 0.95$--$0.97$ for a $19$-CNOT circuit even
though, with multi-restart and $100$ L-BFGS iterations during post-training
refinement, the same topology often reaches $F = 1.00$.

The agent therefore receives a \emph{biased} fidelity signal at training
time. In the regime that should be optimal (low CNOT, K-G shape), the
inner-refined fidelity is lower than the truly attainable fidelity, while
in the slightly-deeper-but-easier regime (parallel CNOTs at $\sim 24$
count) the inner-refined fidelity is genuinely high. The QASER multiplier
$\text{base}^F$ then pulls in favour of the easier-to-fit shape.

This bias is not visible at $n = 2$ because the angle-convergence problem
is trivial there ($\le 3$ rotation parameters between any two CNOTs),
and the inner L-BFGS is well within its sweet spot.

\subsection{Why the post-training refinement does not retrofix this}

Post-training refinement \emph{can} converge tightly: it runs $100$
L-BFGS steps with $16$ random restarts. But it operates on the entries
\emph{already} in the Pareto archive. Since the agent rarely sampled
near-K-G-shaped circuits during training, those circuits never entered
the archive in the first place. Refinement polishes what it is given;
it cannot replace the discrete topology.

The conclusion is unambiguous: \textbf{the bottleneck for matching
Qiskit's CNOT count on $n \ge 3$ is angle-convergence quality at training
time, specifically inside the inner-refinement step}. To fix it we need
either (a) more L-BFGS work per sample, which scales the rollout cost
linearly, or (b) a better optimiser that converges in fewer evaluations.
We argue below that Rotosolve gives us option (b).

% ════════════════════════════════════════════════════════════════════════════
\section{Convergence speed: where time is actually spent}
% ════════════════════════════════════════════════════════════════════════════

A second axis of improvement is the wall-clock time per training run.
Per-circuit averages are $75.5$ s ($n = 2$) and $120.2$ s ($n = 3$). For
a $50$-circuit benchmark this is roughly $1$ hour for $n = 2$ and
$1$h$40$ for $n = 3$ on GPU.

Profiling broadly partitions the per-epoch cost as follows:
\begin{enumerate}[leftmargin=2em]
\item \textbf{Rollout (KV-cached transformer decode).} O(N\_samples $\cdot$ T) — small for the tiny model.
\item \textbf{Inner refinement (12 L-BFGS steps $\times$ 1024 samples).} The dominant cost. Each L-BFGS step
takes one value-and-grad evaluation through the $T_{\max}$-step circuit
\texttt{lax.scan}; with $T_{\max} = 105$ ($n = 3$) this is $\sim 12 \cdot 105 = 1260$ scan iterations per sample.
\item \textbf{Fidelity evaluation for reward.} One forward pass per sample.
\item \textbf{PPO update (4 $\times$ 256-mini-batch passes through the
buffer).} Less expensive than inner refinement because the gradient is
shared across all parameters of the transformer rather than re-traced
per circuit.
\item \textbf{Post-training refinement.} Once at the end:
$|\text{archive}| \cdot 100$ L-BFGS steps $\times 16$ restarts.
Cheap per circuit (small archive size) but quadratic in
$T_{\max}$ at the time of jit compilation.
\end{enumerate}

The two practical levers are (i) reducing the per-step cost of the
inner-loop optimiser, and (ii) reducing the number of inner-loop steps
needed for convergence. L-BFGS scores well on (i) (one cheap evaluation
per step) but poorly on (ii) for highly non-convex parameterised quantum
circuits (Section~\ref{sec:diagnosis}). Rotosolve has the opposite
profile: each sweep is more expensive than an L-BFGS step but a small
constant number of sweeps converges to a global optimum within each
parameter's coordinate. We show next that for our objective Rotosolve's
update is \emph{exact}, not an approximation.

% ════════════════════════════════════════════════════════════════════════════
\section{Rotosolve for hybrid-action GQE}
% ════════════════════════════════════════════════════════════════════════════

\subsection{Why our objective is sinusoidal in every angle}

The process-fidelity cost as a function of one rotation parameter $\theta_i$
is exactly a single Fourier mode plus a constant. Let $g_i(\theta_i) = \cos(\theta_i / 2) I - \E^{i\, 0}\,\sin(\theta_i / 2)\,\sigma_{\alpha_i}$
(an $\mathrm{R}_x$, $\mathrm{R}_y$, or $\mathrm{R}_z$ rotation in the
hardware-aware basis). Then $\Ucirc(\bm\theta) = U_{>i}\, g_i(\theta_i)\, U_{<i}$
with $U_{<i}, U_{>i}$ independent of $\theta_i$, and so
\begin{align*}
\mathrm{tr}\bigl(\Utarget^\dagger\,\Ucirc(\bm\theta)\bigr)
&= \cos(\theta_i / 2)\,\underbrace{\mathrm{tr}\bigl(\Utarget^\dagger\, U_{>i}\, U_{<i}\bigr)}_{=:\,a_i}
- i\,\sin(\theta_i / 2)\,\underbrace{\mathrm{tr}\bigl(\Utarget^\dagger\, U_{>i}\, \sigma_{\alpha_i}\, U_{<i}\bigr)}_{=:\,b_i},
\end{align*}
which after squaring and using $\cos^2(\theta/2) = (1 + \cos\theta)/2$
and $\sin^2(\theta/2) = (1 - \cos\theta)/2$ gives
\begin{equation}
\boxed{\;
F(\theta_i)
\;=\;
\frac{|a_i|^2 + |b_i|^2}{2 d^2}
\;+\;
\underbrace{\frac{|a_i|^2 - |b_i|^2}{2 d^2}}_{=: A_i}\,\cos\theta_i
\;+\;
\underbrace{\frac{\mathrm{Im}(\bar a_i b_i)}{d^2}}_{=: B_i}\,\sin\theta_i.
\;}
\label{eq:single-param-cost}
\end{equation}
Equation~(\ref{eq:single-param-cost}) is the entire content of Rotosolve.
Holding all other parameters fixed, $F(\theta_i)$ is a single sinusoid
plus a constant; three function evaluations are sufficient to reconstruct
it exactly, and its global maximum is a closed-form arctangent. No
gradient, no learning rate, no line search, no restarts \emph{in this
coordinate}.

\subsection{The closed-form update}

Choose three reference shifts $\{0, +\pi/2, -\pi/2\}$ relative to the
current $\theta_i^{(k)}$, evaluate
$F^0_i = F(\theta_i^{(k)})$,
$F^+_i = F(\theta_i^{(k)} + \pi/2)$,
$F^-_i = F(\theta_i^{(k)} - \pi/2)$. Substituting into
Equation~(\ref{eq:single-param-cost}) and absorbing the offset $\theta_i^{(k)}$
yields
\begin{equation}
F\bigl(\theta_i^{(k)} + \alpha\bigr)
\;=\;
A'_i + B'_i\cos\alpha + C'_i\sin\alpha,
\end{equation}
with
\begin{equation}
A'_i = \frac{F^+_i + F^-_i}{2},\quad
B'_i = F^0_i - A'_i,\quad
C'_i = \frac{F^+_i - F^-_i}{2}.
\end{equation}
The exact maximiser of $B'_i\cos\alpha + C'_i\sin\alpha$ is
$\alpha^\star_i = \mathrm{atan2}(C'_i, B'_i)$. The Rotosolve update is
therefore the closed-form
\begin{equation}
\boxed{\;
\theta_i^{(k+1)}
\;=\;
\theta_i^{(k)}
\;+\;
\mathrm{atan2}\!\bigl(F^+_i - F^-_i,\;\;2\,F^0_i - F^+_i - F^-_i\bigr).
\;}
\label{eq:rotosolve-update}
\end{equation}
A \emph{sweep} of Rotosolve cycles through every parametric angle
$\theta_i$ (those whose token is $\mathrm{R}_x, \mathrm{R}_y$, or
$\mathrm{R}_z$) and applies (\ref{eq:rotosolve-update}). Non-parametric
positions ($\mathrm{SX}$, $\mathrm{CNOT}$, $\mathrm{STOP}$) are skipped.

\subsection{Cost accounting and the comparison with L-BFGS}

Let $P$ denote the number of parametric positions in a circuit (typically
$P \in [60, 100]$ for our $n = 3$ runs). One Rotosolve sweep costs $3P$
fidelity evaluations: three per parameter, each a full $T_{\max}$-step
\texttt{lax.scan} forward pass through the same JIT-compiled kernel as
L-BFGS. There is no gradient; the kernel is the existing forward-only
fidelity function.

In comparison, one L-BFGS step costs roughly two evaluations (forward +
\texttt{value\_and\_grad}). Per evaluation L-BFGS is therefore slightly
more expensive (we pay the backward pass), but per \emph{step} L-BFGS
makes a small move; per \emph{sweep} Rotosolve makes the globally optimal
move within each parameter coordinate. Empirically, on parameterised
quantum circuits, the convergence threshold for a fixed final fidelity
behaves as:
\begin{center}
\small
\begin{tabular}{lll}
\toprule
optimiser & evaluations to $F \!\ge\! 0.999$ (typical) & per-evaluation cost \\
\midrule
L-BFGS, single trajectory & $\sim 100$--$300$, often gets stuck         & $\sim 1.5\times$ a forward (grad) \\
L-BFGS with $16$ restarts & $\sim 100\times 16 = 1600$                 & same  \\
Rotosolve, $K$ sweeps     & $3 P K$, e.g.\ $3\!\cdot\! 80\!\cdot\! 4 = 960$ & $\sim 1\times$ a forward \\
\bottomrule
\end{tabular}
\end{center}
For $P = 80$, $K = 4$ sweeps roughly matches the work of $16$ L-BFGS
restarts, with the qualitative advantage that Rotosolve's coordinate
update is \emph{exact} so it cannot get trapped in the per-coordinate
local minima that single-trajectory L-BFGS often does.

\subsection{Where to deploy Rotosolve}

There are two natural deployment points:

\paragraph{(a) Inner refinement (per rollout sample).}
Currently each of the $1024$ rollout samples is polished by $12$
L-BFGS steps at the end of every epoch. Replacing this with $K$ Rotosolve
sweeps gives a stronger fidelity signal to the policy gradient at
training time. The trade-off is per-step cost. With $P \sim 80$ rotation
parameters and $K = 1$ sweep, the inner cost goes from
$\sim 1024 \cdot 12 \cdot 2 \approx 24{,}000$ fidelity evaluations per
epoch to $\sim 1024 \cdot 3 \cdot 80 = 245{,}000$ evaluations per epoch
--- a $10\!\times\!$ slow-down. With $K = 1$ this would be unattractive
unless the per-circuit fidelity convergence is correspondingly tighter.

A more pragmatic option for the inner stage is a \emph{partial} sweep:
update only the $P_{\mathrm{sub}} \le P$ parameters with the largest
expected impact (e.g.\ those whose gates lie on the same qubit as the
nearest CNOT, or simply a random $P_{\mathrm{sub}} = 8$ subset). This
caps the per-sample cost at $24$ fidelity evaluations regardless of the
circuit length while still providing the closed-form coordinate update
for the parameters Rotosolve does touch. We expect this to fix the
biased reward signal that drives the CNOT inflation in
Section~\ref{sec:diagnosis} without paying $10\times$ the wall-clock.

\paragraph{(b) Post-training Pareto refinement.}
This is the most attractive deployment. The archive after training is
small ($\sim 30$--$80$ entries), so per-circuit cost matters less. With
$K = 5$ sweeps and $P \le 100$, post-training refinement costs
$\sim 1500$ fidelity evaluations per circuit, similar to the current
$100\, \text{steps} \times 16\, \text{restarts} = 1600$ L-BFGS evaluations
per circuit. Rotosolve's per-coordinate exactness is expected to give it
a fidelity edge of $0.001$--$0.01$ on the hard cases (the kind of cases
that today plateau just below $F = 0.999$).

\subsection{Limitations and caveats}

Three points to note:

\begin{enumerate}[leftmargin=2em]
\item \textbf{Sequential within a circuit.} Rotosolve must update one
parameter at a time --- the closed-form is exact only conditional on
the other parameters being fixed. We can vmap it across the
\emph{batch} of circuits, but not across the parameters of a single
circuit. L-BFGS, in contrast, updates all parameters simultaneously
each step.

\item \textbf{Saddle points are still possible.} Rotosolve descends
greedily in coordinates and can land at a point where every individual
$\theta_i$ is at its conditional optimum but the joint configuration is
sub-optimal. This is the same failure mode L-BFGS has at saddles, and
the standard cure is the same: random restarts. So multi-restart
remains useful for Rotosolve at post-training, but with fewer restarts
than L-BFGS needs.

\item \textbf{Single-shift parameter only.} The closed form
(\ref{eq:single-param-cost}) holds when $\theta_i$ enters the circuit as
a single $\mathrm{R}_\alpha(\theta_i)$ rotation. Our pool satisfies this
trivially (each parametric token is a single rotation). If we ever add
multi-frequency parametric gates (e.g.\ a continuously parameterised
two-qubit gate), the cost would become a multi-frequency Fourier series
and the basic Rotosolve would need a generalised parameter-shift rule.
For now this is a non-issue.
\end{enumerate}

% ════════════════════════════════════════════════════════════════════════════
\section{Other improvements to consider}
% ════════════════════════════════════════════════════════════════════════════

The Rotosolve discussion above targets the dominant pathology
(angle-convergence bias on $n = 3$). A few smaller knobs are also worth
documenting in case Rotosolve alone does not close the CNOT gap.

\subsection{CNOT-aware archive criteria}

Currently the archive admits any non-dominated $(F, D, C)$ point above
the configured \texttt{fidelity\_floor}. A CNOT-tight circuit at
$F \approx 0.97$ is therefore evicted in favour of a parallel-CNOT
circuit at $F \approx 0.99$, even though the post-training refiner
\emph{would} push the tight circuit's $F$ to $1.00$ if given the chance.

A small change is to add a second admission criterion: if a candidate's
CNOT count is at most $C_{\mathrm{KG}}(n) + \Delta$ for a small slack
$\Delta$ (say $\Delta = 1$), admit it regardless of its current
fidelity (subject to the hard \texttt{fidelity\_floor}). The post-training
refiner then has the chance to polish it to a competitive $F$. This is a
purely archive-side change with no effect on the policy gradient.

\subsection{Curriculum on the inner-refinement budget}

Instead of $12$ L-BFGS steps for every epoch, schedule the inner-step
budget so that early epochs get a small budget (the agent is still
exploring topology and a tight angle fit is wasted) and late epochs get a
large budget (the agent has converged on a topology and we want a clean
fidelity signal). This is a small change to \texttt{config.yml} and the
trainer, and it spends the same total wall-clock more efficiently.

\subsection{Increase $w_C$ (CNOT weight) only late in training}

The current $w_C = 2.0$ is a compromise between not over-pressuring the
agent early (when it doesn't yet have a good topology) and pushing it
toward CNOT compactness late. A simple curriculum
$w_C(t) = w_C^{\text{init}} + (w_C^{\text{final}} - w_C^{\text{init}}) \cdot t / T_{\max}$
with $w_C^{\text{init}} = 1$ and $w_C^{\text{final}} = 3$ would let the
agent first find any high-$F$ topology and then prune CNOTs from it.

\subsection{Initial-angle restarts at the inner stage}

Even without changing the optimiser, doing two L-BFGS restarts per inner
sample (one from the RL angles, one from $\mathrm{Uniform}[-\pi, \pi]$)
would catch the most common single-trajectory failures, at $2\times$
inner cost. This is a cheap experiment to disentangle ``L-BFGS is
intrinsically inadequate'' from ``L-BFGS gets stuck from the RL
initialisation specifically''.

% ════════════════════════════════════════════════════════════════════════════
\section{Proposed configuration surface}
% ════════════════════════════════════════════════════════════════════════════

If we adopt Rotosolve as a configurable optimiser the simplest exposure
in \texttt{config.yml} is one knob per refinement stage, alongside a
small number of method-specific hyperparameters:

\begin{lstlisting}[language=yaml]
policy:
  inner_refine:
    method: "lbfgs"          # "lbfgs" | "rotosolve" | "none"
    steps: 12                # L-BFGS step count (when method=lbfgs)
    sweeps: 1                # Rotosolve sweeps   (when method=rotosolve)
    sweep_subset: null       # Rotosolve only: int -> partial-sweep size,
                             #                 null -> full sweep over P
    lr: 0.1                  # L-BFGS learning rate

refinement:
  method: "lbfgs"            # "lbfgs" | "rotosolve"
  steps: 100                 # L-BFGS step count
  sweeps: 5                  # Rotosolve sweeps
  num_restarts: 16           # used by both methods
  apply_simplify: true
  lr: 0.1
\end{lstlisting}

The \texttt{sweep\_subset} knob is the inner-loop pragmatic compromise
(Section~5d): if set to e.g.\ $8$, each inner Rotosolve call updates only
the $8$ parameters whose adjacent gates have the largest joint two-qubit
``footprint''. This keeps the inner-stage cost at the same magnitude as
the current L-BFGS while still benefiting from Rotosolve's exactness.

\subsection{Backwards compatibility}

Setting \texttt{method: "lbfgs"} everywhere reproduces today's behaviour
exactly. The new \texttt{rotosolve} branch is purely additive: it shares
the existing JIT-compiled fidelity kernel, requires no changes to the
PPO loss, the rollout function, or the Pareto archive, and does not
alter the buffer schema.

\subsection{Implementation surface}

\begin{itemize}[leftmargin=2em]
\item \textbf{New module: \texttt{rotosolve.py}.} A single function
\texttt{rotosolve\_refine\_batch(token\_ids, init\_angles, evaluator,
sweeps, parametric\_mask, sweep\_subset=None) -> (refined\_angles,
refined\_fidelities)} implementing the loop in
Equation~(\ref{eq:rotosolve-update}). Internals are pure JAX with
\texttt{jax.lax.fori\_loop} over sweeps and \texttt{jax.vmap} over the
batch of circuits.

\item \textbf{\texttt{AngleRefiner} dispatch.} The existing class gains
a constructor switch on \texttt{method}; for \texttt{rotosolve} it dispatches
to the new module's batched function rather than the existing L-BFGS
\texttt{fori\_loop}.

\item \textbf{\texttt{config.py}.} Three new fields under
\texttt{policy.inner\_refine.\{method, sweeps, sweep\_subset\}} and
two new fields under \texttt{refinement.\{method, sweeps\}}. The
existing \texttt{steps} and \texttt{lr} fields stay.

\item \textbf{Tests.} A unit test that on a known target,
\texttt{Rotosolve} reaches $F \ge 0.999$ in $K \le 4$ sweeps, and a
property test that one Rotosolve sweep does not \emph{decrease}
fidelity (it should be monotone non-decreasing by construction).
\end{itemize}

The implementation is small ($\sim 80$ lines for the optimiser, $\sim 40$
for the dispatch, $\sim 30$ for tests) and the risk surface is
self-contained: if Rotosolve underperforms in practice we can always
revert to \texttt{method: "lbfgs"} per stage without touching anything
else.

% ════════════════════════════════════════════════════════════════════════════
\section{Expected outcomes}
% ════════════════════════════════════════════════════════════════════════════

Based on the diagnosis in Section~\ref{sec:diagnosis} and the
parameter-counting in Section~4, our expectations after this change are:

\begin{itemize}[leftmargin=2em]
\item \textbf{$n = 2$.} No measurable change. The angle convergence is
already trivial here; both L-BFGS and Rotosolve hit $F = 1$ within a few
iterations. The benchmark should remain a clear win on every axis.

\item \textbf{$n = 3$, fidelity.} A small absolute gain, $\bar F$ from
$0.99977$ to $\ge 0.99990$, driven mostly by post-training refinement
no longer leaving entries stuck just below $F = 0.999$.

\item \textbf{$n = 3$, CNOT count.} The intended outcome: median
$\Delta C$ from $+5$ down to $\le +2$, with at least a third of
instances matching Qiskit's $19$-CNOT count and a small but non-zero
fraction beating it (the Khaneja-Glaser optimum is $14$ for $n = 3$ and
some of the $50$ targets should be expressible at exactly that count).

\item \textbf{$n = 3$, depth.} Slight regression expected (median
$\Delta D$ from $-5$ to $-3$): once the agent is no longer rewarded for
parallel-CNOT redundancy, the circuits become slightly less wide.
Total gate count is expected to shift only marginally because rotations
are not the bottleneck.

\item \textbf{Wall-clock per circuit.} Inner stage with
\texttt{sweep\_subset = 8}: same magnitude as today. Post-training
stage with \texttt{sweeps = 5}: same magnitude as today. The extra
implementation is essentially free in wall-clock at the same fidelity
target.

\item \textbf{$n = 4$.} The angle-convergence problem is more severe at
$n = 4$ (we observed an $F \approx 0.98$ plateau in earlier runs on
this benchmark). Rotosolve at the inner stage in particular is more
likely to lift the plateau here than at $n = 3$, because the relative
gap between L-BFGS and Rotosolve grows with parameter count. We
\emph{do not} expect Rotosolve alone to fix the $n = 4$ plateau --- the
deeper bottleneck there is the tiny model's discrete capacity --- but
we expect a meaningful contribution.
\end{itemize}

% ════════════════════════════════════════════════════════════════════════════
\section{Summary}
% ════════════════════════════════════════════════════════════════════════════

The benchmark on $50$ brickwork-Haar targets at $n = 2$ and $n = 3$
shows the pipeline produces high-fidelity, structurally-near-optimal
circuits at $n = 2$, and matches or exceeds Qiskit on fidelity, depth,
and total gate count at $n = 3$, but loses on CNOT count by a median
of $+5$. The CNOT delta is a symptom of biased fidelity signalling at
training time: the inner-refinement L-BFGS does not converge tightly
enough on near-Khaneja-Glaser-optimal topologies, so the agent settles
on slightly larger but easier-to-fit ones.

We propose adding Rotosolve as a configurable optimiser at both
refinement stages. For our process-fidelity objective, Rotosolve's
update is exact in each parameter coordinate (Equation
\ref{eq:rotosolve-update}), requires no learning rate, and converges in
a small number of full sweeps. The implementation surface is small and
fully backward-compatible: setting the method to \texttt{lbfgs}
recovers today's pipeline.

\paragraph{References.}
\begin{itemize}[leftmargin=2em]
\item Ostaszewski, Grant, Benedetti.\ \emph{Structure optimization for parameterized quantum circuits.}\ \texttt{arXiv:1905.09692} (Rotosolve).
\item Wierichs, Izaac, Wang, Lin.\ \emph{General parameter-shift rules for quantum gradients.}\ \texttt{Quantum 6, 677 (2022)}.
\item Niu et al.\ \emph{Hybrid Action Reinforcement Learning for Quantum Architecture Search.}\ \texttt{arXiv:2511.04967} (HyRLQAS).
\item Moflic et al.\ \emph{QASER: Breaking the Depth vs.\ Accuracy Trade-Off for Quantum Architecture Search.}\ \texttt{arXiv:2511.16272}.
\item Shende, Markov, Bullock.\ \emph{Minimal universal two-qubit CNOT-based circuits.}\ \texttt{arXiv:quant-ph/0308033}.
\end{itemize}

\end{document}
